The following astronomical data are given.
\begin{itemize}
\item The masses of Mars $M_M$ and the Sun $M_\odot$ are expressed relative to Earth's mass $M_\oplus$, with $M_M=0.107M_\oplus$ and $M_\odot=333054M_\oplus$.
\item The mean Mars--Sun distance is $r_{MS}=1.524\,\mathrm{au}$ and Mars has a circular orbit.
\item The mean Earth--Sun distance is $r_{ES}=1\,\mathrm{au}=1.496\times10^{11}\,\mathrm m$ and Earth has a circular orbit.
\item Mars's moon Phobos has mass $M_{\mathrm{Phobos}}=1.5\times10^{-8}M_M$.
\item Earth's sidereal period is $P_{\mathrm{Sid}\oplus}=1$ year $=365.2564$ days.
\item Phobos has average radius $R_{\mathrm{Phobos}}=12.5\,\mathrm{km}$ and orbital radius $d_{\mathrm{Phobos}M}=9380\,\mathrm{km}$.
\item Phobos is spiraling inward towards Mars at the constant rate of 1.8\,cm/yr.
\item You may use $(1+x)^n\approx1+nx$ when $x\ll1$.
\item The Roche limit is the boundary inside which the tidal force of the planet on a small mass is more than its self-gravity.
\end{itemize}
\begin{description}
\item[(a)] Derive the formula for the L1 Lagrange point (distance between secondary body and the L1 point) and calculate its numerical value for the Sun--Mars system. Find the ratio of Phobos's orbital radius to this distance.
\item[(b)] Derive the formula for the L2 Lagrange point (distance between secondary body and the L2 point) and calculate its numerical value for the Mars--Phobos system.
\item[(c)] Derive the formula for the Roche limit of the Mars--Phobos system. How long would it take Phobos to reach this limit? It is known that it approaches the surface of the planet during its rotation and will either collide with it or begin to disintegrate, forming a ring. The satellite is considered to be solid.
\item[(d)] Calculate the distance $l$ between the Lagrange points L4 and L5 for the Sun--Mars system.
\end{description}
